Es sei $\gamma$ ein Halbweg in $M$ wie angegeben und sei

\mavergleichskette
{\vergleichskette
{ P
}
{ \in }{  U
}
{ \subseteq }{ M
}
{    }{
}
{   }{
}
}
{}{}{}
ein Kartengebiet mit einer Karte
\maabbdisp {\alpha} { U } {V \subseteq  H_{\leq 0}
} {}
mit

\mavergleichskettedisp
{\vergleichskette
{  H_{\leq 0}
}
{ =} { \left\{ (x_1,x_2 , \ldots, x_n)  \mid   x_1 \leq 0        \right\}
}
{ } {
}
{ } {
}
{ } {
}
}
{}{}{}
und mit

\mavergleichskettedisp
{\vergleichskette
{ Q
}
{ =} { \alpha(P)
}
{ =} { \left(  0   , \,  a_2  , \,   \ldots , \,  a_n                \right)
}
{ } {
}
{ } {
}
}
{}{}{.}
Unter der Tangentialabbildung wird

\mavergleichskette
{\vergleichskette
{ v
}
{ = }{ \gamma'(0)
}
{  }{
}
{    }{
}
{   }{
}
}
{}{}{}
auf den Ableitungsvektor von

\mavergleichskette
{\vergleichskette
{ \delta
}
{ = }{ \alpha \circ \gamma
}
{  }{
}
{    }{
}
{   }{
}
}
{}{}{}
im Nullpunkt abgebildet. Da $\gamma$ in $M$ verläuft, verläuft $\delta$ ganz in $H_{\leq 0}$. Daher ist im Differenzenquotient

\mathdisp {{ \frac{   \left(  \delta_1(t)   , \,   \delta_2(t)  , \,   \ldots , \,   \delta_n(t)                \right)  }{ t } }} {  }

sowohl $t$ negativ als auch $\delta_1(t)$ negativ und daher ist

\mavergleichskettedisp
{\vergleichskette
{  \delta_1'(0)
}
{ \geq} { 0
}
{ } {
}
{ } {
}
{ } {
}
}
{}{}{,}
gehört also zur positiven Hälfte.

Es gehöre umgekehrt
\mathl{T_P(\alpha)(v)}{} zur positiven Hälfte. Dann verläuft die Gerade
\maabbeledisp {\delta} {\R} {\R^n
} { t } { Q+tv
} {}
für

\mavergleichskette
{\vergleichskette
{ t
}
{ \leq }{  0
}
{  }{
}
{    }{
}
{   }{
}
}
{}{}{}
ganz in der negativen Hälfte. Der Weg
\mathl{\alpha^{-1} \circ \delta}{} ist auf einem negativen Intervall definiert und landet in $M$ und ist eine Halbwegrealisierung von $v$.

 [[Kategorie:Latexseite]]
